\documentclass{article}
\usepackage[T1]{fontenc}
\usepackage{lmodern}
\usepackage{graphicx}
\usepackage{amsmath,amssymb}
\usepackage[a4paper,margin=2cm]{geometry}
\usepackage[absolute,overlay]{textpos}
\setlength{\TPHorizModule}{1mm}
\setlength{\TPVertModule}{1mm}
\usepackage[indonesian]{babel}
\usepackage{hyperref}
\setlength{\emergencystretch}{2em}
% unit: Halaman 92

\begin{document}

% === Halaman 92 (llm) ===

\begin{itemize}
    \item Jika pasangan plainteks dan cipherteks tidak tersedia, maka Affine Cipher masih bisa diserang dengan \textit{exhaustive key search}.
    \item sebab ada 25 pilihan untuk $b$ dan 12 buah nilai $m$ yang relatif prima dengan 26 (yaitu 1, 3, 5, 7, 9, 11, 15, 17, 19, 21, 23, dan 25).
    \item Dengan mencoba sebanyak $25 \times 12 = 300$ susunan $b$ dan $m$, maka nilai $b$ dan $m$ yang benar dapat ditemukan dengan mudah, karena 300 susunan $b$ dan $m$ adalah jumlah yang sangat sedikit.
\end{itemize}

\end{document}
